% !TEX root = ../main.tex
% =====================================================================
%  §6  General Form Of The Higgs Branch Contribution   （原 PDF pp. 58--59）
%  §7  The Coulomb Branch Contribution aka The u-Plane integral（原 PDF pp. 60--66）
%  公式编号 (6.1)--(6.7)、(7.1)--(7.24)，逐条保留原文编号；图位 3 处。
% =====================================================================

\section{希格斯分支贡献的一般形式}

\begin{keytakeaways}
\item \textbf{物理图像：} 希格斯分支的贡献只来自两个阿贝尔真空 $u=\pm\Lambda^{2}$；一旦把通量扇区 $\lambda$ 固定，低能有效理论就是“阿贝尔矢量多重态 $+$ 磁单极超多重态”的标准形式，其路径积分恰好局域化到塞伯格--威滕方程的解空间上。
\item \textbf{数学工具：} 模空间 $\cM(\lambda)$ 的虚拟维数公式 (6.3)（由脚注 31，$2\chi+3\sigma=p_1(\Lambda^{+})$，故维数必为偶数），以及把 $u$ 按鬼数为 $2$ 的场 $a_d$ 展开、再用留数把 $\cM(\lambda)$ 上的积分约化为 SW 不变量的技巧。
\item \textbf{核心结论：} 式 (6.1)--(6.7) 把 $Z_{IR,Higgs}$ 写成 SW 不变量乘上四个未知耦合 $T_h,C,P_h,E_h$ 的函数；剩下的任务只是定出这些耦合，而这一步由第 8 节借助库仑分支的穿墙数据完成。
\end{keytakeaways}

现在概述如何用未知函数 $T_h,C,P,E$ 来计算 $Z_{IR,Higgs}$。

如前所述，$Z_{IR,Higgs}$ 是 $u=\Lambda^{2}$ 与 $u=-\Lambda^{2}$ 两项之和。$u=-\Lambda^{2}$ 处的贡献可由第 5.4.2 节所述对称性给出：该对称性取 $u\to-u$、$a\to\ii a$。因此我们只需关注 $u=\Lambda^{2}$ 这一项，它可写为
% [Tutor Note: 这里的 $\ZZ_2$ 对称性（$u\to-u$、$a\to\ii a$）把双荷子点 $u=-\Lambda^{2}$ 的贡献与磁单极点 $u=\Lambda^{2}$ 的贡献联系起来；注意 $\dd u/\dd a$ 在该变换下也带相位，故两个求和项之间的相对相位不是平凡的 $1$，可对照 (9.1) 中第二项多出的 $\ii^{\chi_h}$。]
\begin{equation}
Z_{IR,Higgs,\Lambda^{2}}
=
\sum_{\lambda\in\Gamma_w}
\left\langle
\exp\left(2pu+\frac{\ii}{4\pi}\int_{\Sigma}\frac{\dd u}{\dd a_d}F(A_d)+\Sigma^{2}T_h(u)\right)
\right\rangle_{u=\Lambda^{2},\lambda}
\tag{6.1}
\end{equation}
其中我们对自旋$^{c}$结构的第一陈类 $\lambda$ 求和，如式 (4.75) 所示。在写 $2$-可观测量时我们作了少许简化，即在此假定 $b_1=0$，从而可以丢掉含 $\psi^{2}$ 的项。

现在我们在固定的“通量扇区”$\lambda$ 中计算路径积分。与轻的“磁单极超多重态”耦合的低能有效作用量取标准的 Mathai--Quillen（MQ）形式，正好可以局域化到 SW 方程上。因此在固定通量扇区中路径积分为
% [Tutor Note: 固定通量扇区之后，路径积分为什么就等于模空间的交截数：局域化把测度压到 $Q$-不动点集上，即 SW 方程的解空间 $\cM(\lambda)$；而积分号内除 $e^{2pu}$ 与 $\Sigma^{2}T_h$ 之外的因子都是 $\cM(\lambda)$ 上特定次数的上同调类（$C^{\lambda^2}$ 来自 2-形式的插入，$P^{\sigma}$ 来自 $\sigma$ 个点的插入，$E^{\chi}$ 来自与欧拉示性数耦合的因子），故路径积分正是这些类的交截数。式 (6.2) 没有写出明显的测度并非笔误：模空间的维数已被 (6.3) 固定为 $2n(\lambda)$，积分只有在次数相配时才非零，这也就是 (6.6) 定义的 SW 不变量只取整数值的原因。]
\begin{equation}
\int_{\cM(\lambda)}
e^{2\pi\ii(\lambda_0^2+\lambda\cdot\lambda_0)}
e^{\,2pu+\ii\frac{\dd u}{\dd a_d}\Sigma\cdot\lambda+\Sigma^{2}T_h(u)}
\,C(u)^{\lambda^{2}}P(u)^{\sigma}E(u)^{\chi}.
\tag{6.2}
\end{equation}
这里 $\cM(\lambda)$ 是基于自旋$^{c}$结构 $\lambda$ 的塞伯格--威滕方程解的模空间。已知它是光滑、紧致、可定向的，其维数为\footnote{量 $2\chi+3\sigma$ 在 SW 理论中出现的频率极高。它是自对偶 2-形式实丛 $\Lambda^{+}$ 的第一庞特里亚金类。对任何复结构都存在典范类，且 $K^{2}=2\chi+3\sigma$。}
\begin{equation}
\vdim\cM(\lambda)=\frac{(2\lambda)^{2}-(2\chi+3\sigma)}{4}:=2n(\lambda),
\qquad n(\lambda)\in\ZZ
\tag{6.3}
\end{equation}
% [Tutor Note: $(2\lambda)^{2}$ 指相交形式下的自交数；$\vdim$ 必为偶数这一点直接来自脚注 31：$2\chi+3\sigma=p_1(\Lambda^{+})$ 与 $2\chi+3\sigma$ 同奇偶，故分子被 $4$ 整除。原文断言 $\cM(\lambda)$ 光滑紧致可定向，这只是“一般度规、无阿贝尔瞬子”时的结论；当 $b_2^{+}=1$ 且存在可约解（$M=0$、$F^{+}=0$）时模空间退化，这正是第 8 节穿墙的来源。]
模空间 $\cM(\lambda)$ 是如下无穷维流形的子流形：
\begin{equation}
\bigl(\cA\times\Gamma(S^{+}\otimes L)\bigr)/\mathcal{G}
\tag{6.4}
\end{equation}
其中 $\cA$ 现在是自旋$^{c}$联络的空间：它是以 $X$ 上整体定义的纯虚 $1$-形式空间为挠子的仿射空间。规范变换群就是 $\mathcal{G}=\mathrm{Map}(X,U(1))$。它作用在自旋$^{c}$联络上为 $\nabla\to\nabla+u^{-1}\dd u$，作用在磁单极场上为 $M\to uM$。
% [Tutor Note: 〔关于 (6.4) 与锥顶的补充〕$\Gamma(S^{+}\otimes L)$ 是无穷维复希尔伯特空间，整体 $U(1)$ 变换以 $M\mapsto\ee^{\ii\theta}M$ 作用，故其商是复射影空间 $\CC\PP^{\infty}$ 上的锥，锥顶即 $M=0$。原文“该锥的尖端只有在存在阿贝尔瞬子 $F^{+}=0$ 时才与 $\cM(\lambda)$ 相交”一句说得极简，其含义是：$M=0$ 时 SW 方程退化为 $F^{+}=0$，即线丛 $L$ 上的阿贝尔瞬子；只有此时锥顶（它本身不是流形，且其邻域是奇异的）才落在模空间里，$\cM(\lambda)$ 才可能出现额外玻色零模、非紧端与 $b_2^{+}=1$ 时的穿墙（对照第 7.3 节 (7.16) 之后关于 $\lambda_+=0$ 的讨论）。若不存在阿贝尔瞬子，锥顶与 $\cM(\lambda)$ 不相交，$\cM(\lambda)$ 完全落在“远离锥顶”的光滑部分，可以安全截断。又 $\CC\PP^{\infty}$ 的上同调环为 $\ZZ[x]$、$x$ 为 $2$ 次类，这正是 $a_d$ 的几何来源。]
在通常的阿贝尔规范理论中，$U(1)$ 联络的同构类 $\cA/\mathcal{G}$ 就是微分上同调群 $\check{H}^{2}(X)$。它是一个无穷维阿贝尔群，可以（非典范地）写成 $T^{b_1}\times V\times H^{2}(X;\ZZ)$，其中 $T^{b_1}$ 是调和 $1$-形式模掉具有整周期的那些所得的环面，$H^{2}(X;\ZZ)$ 给出第一陈类，$V$ 是无穷维向量空间。在目前情形下我们处理的是自旋$^{c}$联络，但空间几乎相同：自旋$^{c}$联络的空间是以 $\Omega^{1}(X)$ 为模型的仿射空间，其商空间为 $T^{b_1}\times V$，$V$ 为无穷维向量空间。整体 $U(1)$ 规范变换不作用于自旋$^{c}$联络，但确实作用于旋量场，因此在拓扑上 $(\nabla,M)$ 的规范不等价位形空间可写为 $T^{b_1}\times V\times\Gamma(S^{+}\otimes L)/U(1)$，其中最后一个因子是旋量截面模掉整体规范变换所得的空间。这是一个无穷维希尔伯特空间模掉一个 $U(1)$，因此是 $\CC\PP^{\infty}$ 上的锥。该锥的尖端只有在存在阿贝尔瞬子 $F^{+}=0$ 时才与模空间 $\cM(\lambda)$ 相交。否则 $\CC\PP^{\infty}$ 的上同调环就是由 $2$ 次类 $x$ 生成的 $\ZZ[x]$。
% [Tutor Note: 原文此段用 $u$ 同时表示 $U(1)$ 规范变换的函数（$\nabla\mapsto\nabla+u^{-1}\dd u$、$M\mapsto uM$）与库仑分支的真空期望值 $u=\langle\tr\phi^{2}\rangle$，二者毫无关系，纯属沿用了同一字母。阅读时勿把规范变换的 $u$ 与 $u$ 平面上的坐标混淆。]

以上这些论断在教科书 \cite{ref17,ref48,ref42} 中均有证明。

在拓扑场论中，上同调类 $x$——在差一个归一化的意义下，并再次援引对应关系 (3.9)——就是鬼数为 $2$ 的场 $a_d$。\footnote{与因子 $T^{b_1}$ 相伴的上同调类，当然是 $1$-可观测量的 IR 版本，后者可以在非单连通四维流形上定义。} 我们对积分 (6.2) 的解释应当是作展开
\begin{equation}
u=\Lambda^{2}-2\ii\Lambda a_d+\OO(a_d^{2})
\tag{6.5}
\end{equation}
这一展开作为椭圆函数的表达式是完全已知的。于是我们定义 \emph{塞伯格--威滕不变量}为
\begin{equation}
SW(\lambda):=\int_{\cM(\lambda)}a_d^{n(\lambda)}
\tag{6.6}
\end{equation}
它是一个整数，因此自旋$^{c}$结构 $\lambda$ 对 $Z_{IR,Higgs,\Lambda^{2}}$ 的贡献可表为
% [Tutor Note: (6.7) 是“只保留 $a_d$ 的最低次项”的留数公式：$\cM(\lambda)$ 的维数是 $2n(\lambda)$，故 $\int_{\cM(\lambda)}a_d^{n(\lambda)}$ 之外，只有展开式 $(e^{2pu+\cdots}C^{\lambda^{2}}P_h^{\sigma}E_h^{\chi})$ 中含 $a_d^{n(\lambda)}$ 的项能在留数中存活，其余项给出 $a_d$ 的非负幂、留数为零。展开式中用省略号隐去的正是这些项（含 $\Sigma\cdot\lambda$ 的线性项等）。这一手段合法，前提是 (6.5) 在 $a_d=0$ 邻域解析，而这由椭圆函数的显式表达式保证。]
\begin{equation}
SW(\lambda)\,
\mathrm{Res}_{a_d=0}\frac{\dd a_d}{a_d^{1+n(\lambda)}}
\left(e^{\,2pu+\cdots}C^{\lambda^{2}}P_h^{\sigma}E_h^{\chi}\right)
\tag{6.7}
\end{equation}
在不对 $T_h,C,P_h,E_h$ 有显式了解的情况下，我们只能做到这一步。

\section{库仑分支贡献，即 $u$ 平面积分}

\begin{keytakeaways}
\item \textbf{物理图像：} $b_2^{+}>0$ 时标度论证使 $u$ 平面积分退化为树图级的有限维积分；费米子零模的计数（$\eta\chi$ 的配对）强制 $b_2^{+}=1$，否则整个库仑分支贡献为零。
\item \textbf{数学工具：} 把 $a$ 平面上的测度改写为 $\tau$ 平面上的模形式测度（$\Theta$ 本质上是 theta 函数，$H(\tau)$ 只依赖 $\tau$ 不依赖 $\bar\tau$），用 $1/y$ 展开与三个尖点处的 $q$ 展开控制积分。
\item \textbf{核心结论：} 得到显式积分 (7.6)--(7.7)；其度规变分化为边界围道积分 (7.19)，在各墙 $W(\lambda)=\{\omega:\omega\cdot\lambda=0\}$ 上呈 $\delta$ 函数，穿墙由模形式的傅里叶系数给出；取 $\alpha^{\chi}\beta^{\sigma}=2^{(2+3\sigma)/4}/\pi$ 即精确复现格策 (格策) 关于唐纳森多项式跳变的公式。
\end{keytakeaways}

现在我们来计算 $Z_{IR,Coul}$。

上文已经描述了进行 $u$ 平面积分所需的全部要素。一个重要的标度论证 \cite{ref38} 表明：当 $X$ 有 $b_2^{+}>0$ 时，结果由树图级路径积分决定——可以忘掉单圈行列式与非平凡的费曼图。
% [Tutor Note: “标度论证”的具体内容是：把度规整体缩放 $g\to t^{2}g$，作用量中依赖度规的部分按 $t$ 的幂次分层；$b_2^{+}>0$ 时可以借助 $Q$-恰当项把单圈行列式与非平凡费曼图的贡献写成随 $t\to0$ 消失的全导数，故只剩树图项。原文紧随其后的全大写提醒（见下）指出 $b_2^{+}=0$ 时该论证失效、只有单圈项存活，故 $b_2^{+}=0$ 的流形（如 $S^{4}$）须另作处理。]
于是我们立刻只剩下一个有限维积分：
\begin{equation}
\int(\dd a\,\dd\bar a)(\dd\eta\,\dd\chi\,\dd\psi)(\dd A\,\dd D)\,
e^{S+\Delta_{\mathrm{grv}}S}\,
e^{\,2pu+\OO_{IR,c}(\Sigma)+\Sigma^{2}T_c(u)}
\tag{7.1}
\end{equation}

\begin{center}
\begin{minipage}{0.92\linewidth}
\small 关于 $b_2^{+}=0$ 的说明：此时只有单圈图有贡献，这一情形可能值得研究。\footnote{〔原稿边注〕原文此处以全大写排印：“REMARK ON $b_2^{+}=0$. JUST ONE-LOOP CONTRIBUTES AND COULD BE INTERESTING.”，系作者自注，非正文。}
\end{minipage}
\end{center}

\subsection{费米子积分}

先考虑费米子积分：
\begin{enumerate}
\item $\eta$ 的零模空间 $H^{0}(X;\RR)$ 是一维的。
\item $\psi$ 的零模空间 $H^{1}(X;\RR)$ 是 $b_1$ 维的。
\item $\chi$ 的零模空间 $H^{2,+}(X;\RR)$ 是 $b_2^{+}$ 维的。
\end{enumerate}

关于费米子积分有三点说明：

\begin{enumerate}
\item 注意，定义格拉斯曼积分 $\dd\eta\,\dd\psi\,\dd\chi$ 需要选取一个定向。与对 $w_2(P)$ 的提升 $2\lambda_0$ 的选取所带来的依赖一样，这个定向问题在唐纳森理论中有漂亮的对应：$H^{0}\oplus H^{1}\oplus H^{2,+}$ 的定向选取决定了瞬子模空间 $\cM$ 的定向。参见 \cite{ref8} 命题 7.1.39，第 282 页。看到这些“精细结构细节”在物理进路中得到镜像，令人欣慰。
\item 为简单起见，我们假定 $X$ 单连通。这意味着可以丢掉 $\psi$ 积分，而且可以丢掉作用量 (5.35) 以及各可观测量中 $\psi$ 的依赖。这使方程大为简化。$b_1\neq0$ 时的方程已在 \cite{ref38} 与 \cite{ref30} 中算出。
\item 此时看一下作用量 (5.35) 即可发现，$\chi$ 的零模总是与 $\eta$ 的零模以组合 $\eta\chi$ 的形式一起出现。由于 $\eta$ 的零模空间只有一维，要使 $u$ 平面积分非零，就必须有 $b_2^{+}(X)=1$。这或许令人气馁，但我们仍要继续。注意，把 $\eta$ 与 $\chi$ 积分掉会带来一个因子
\begin{equation}
\frac{\dd\bar\tau}{\dd a}(D+F)_{+}
\tag{7.2}
\end{equation}
\end{enumerate}
% [Tutor Note: (7.2) 是本节的关键约束。$\chi$ 是自对偶 $2$-形式、鬼数 $-1$；$\eta$ 是标量鬼。作用量中二者只以双线性 $\eta\chi$ 出现，而格拉斯曼积分要非零必须让每个零模恰好被“吃掉”一次：$\eta$ 只有一维零模，故只能容纳一个 $\chi$ 零模，于是必须有 $b_2^{+}=1$。$b_2^{+}>1$ 时 $\eta\chi$ 型双线性张成的维数不匹配，$u$ 平面积分整体为零——这一点在 §9 中正是“$b_2^{+}>1$ 时 $Z_{IR,Coul}=0$”的来源。$b_2^{+}=1$ 时 $\chi$ 零模可写为 $\chi=\chi_+\omega$，剩下的因子 $\frac{\dd\bar\tau}{\dd a}(D+F)_+$ 与测度 $\dd a\,\dd\bar a$ 合并，正好给出 (7.6) 中的 $\dd a\,\dd\bar a\,\frac{\dd\bar\tau}{\dd\bar a}$。]

此时有一个有用的观察：当 $b_2^{+}=1$ 时，上同调空间 $H^{2}(X;\RR)$ 带有洛伦兹型度规。一旦选定了 $H^{2,+}$ 的定向，就存在唯一的自对偶类 $\omega$ 使 $\omega\cdot\omega=1$，于是 $\omega$ 落在“质壳”上。

\begin{figure}[H]
  \centering
  \figplaceholder{〔原文此处（印刷页第 60 页，PDF 第 61 页）留有图位“FIGURE: HYPERBOLOID”，PDF 中未插入实际图形〕}
  \caption{原文图题“HYPERBOLOID”：$b_2^{+}=1$ 时满足 $\omega\cdot\omega=1$ 的自对偶类 $\omega$ 在 $H^{2}(X;\RR)$ 中描出的双曲面（即“质壳”）。}
  \label{fig:hyperboloid}
\end{figure}
% [原 PDF 第 61 页]

因此，当 $b_2^{+}=1$ 时我们采用如下约定：$2$-形式可写为
\begin{equation}
F=\omega F_{+}+F^{-}
\tag{7.3}
\end{equation}
其中 $F_{+}$ 是一个标量。这一点在写 (7.2) 时已被隐含使用。

\subsection{光子路径积分}

对规范场积分是直截了当的。如前所述，我们对以 $\lambda\in\Gamma_{\xi}$ 标记的通量扇区求和。结果是作替换
\begin{equation}
F\to4\pi\lambda
\tag{7.4}
\end{equation}
并得到
\begin{equation}
\frac{e^{2\pi\ii\lambda_0^{2}}}{\sqrt{\operatorname{Im}\tau}}
\sum_{\lambda\in\Gamma_{\xi}}
e^{-\ii\pi\bar\tau\lambda_+^{2}-\ii\pi\tau\lambda_-^{2}-\ii\frac{\dd u}{\dd a}(\Sigma,\lambda_-)}
(-1)^{(\lambda-\lambda_0)\cdot w_2(X)}
\left[4\pi\lambda_+ + D\right]
\tag{7.5}
\end{equation}
% [Tutor Note: (7.5) 中 $(-1)^{(\lambda-\lambda_0)\cdot w_2(X)}$ 是自旋$^{c}$ 结构的分级因子，它依赖于 $w_2(P)$ 的提升 $2\lambda_0$ 的选取；换提升 $\lambda_0\to\lambda_0+\mu$（$\mu$ 为整上同调类）时该因子改变 $(-1)^{\mu\cdot w_2(X)}$，须与 $e^{2\pi\ii\lambda_0^{2}}$、$e^{2\pi\ii\lambda\cdot\lambda_0}$ 的相应相位联合考虑，才能保证结果与提升无关（此处用到 $w_2(X)$ 是 $H^{2}(X;\ZZ)$ 的特征向量）。此外注意 $\lambda_{\pm}$ 是对同一整类 $\lambda$ 在 $H^{2,\pm}$ 上的投影，“通量”是 $\lambda$ 而非 $\lambda_\pm$。]

\subsection{$u$ 平面积分的最终表达式}

最后我们只需对辅助场 $D$ 的零模作高斯积分。这是一个一维积分，最终表达式为 $Z_{IR,Coul}=Z_u^{\xi}(p,\Sigma)$，其中
\begin{equation}
Z_u^{\xi}(p,\Sigma)
=
\alpha^{\chi}\beta^{\sigma}
\int
\dd a\,\dd\bar a\,
\frac{\dd\bar\tau}{\dd\bar a}
\left(\frac{\dd u}{\dd a}\right)^{\chi/2}
\Delta^{\sigma/8}
e^{\,2pu+\Sigma^{2}T_c(u)}
\Theta
\tag{7.6}
\end{equation}
\begin{equation}
\Theta=
\frac{e^{\frac{\Sigma_+^{2}}{8\pi y}\left(\frac{\dd u}{\dd a}\right)^{2}}}{\sqrt{y}}
\,e^{2\pi\ii\lambda_0^{2}}
\sum_{\lambda\in\Gamma_{\xi}}
e^{-\ii\pi\bar\tau\lambda_+^{2}-\ii\pi\tau\lambda_-^{2}-\ii\frac{\dd u}{\dd a}(\Sigma,\lambda_-)}
(-1)^{(\lambda-\lambda_0)\cdot w_2(X)}
\left[\lambda_+ + \frac{\ii}{4\pi y}\frac{\dd u}{\dd a}\Sigma_+\right]
\tag{7.7}
\end{equation}
其中我们定义 $y=\operatorname{Im}\tau$，而显含 $y$ 的项来自对 $D$ 的高斯积分。\footnote{若先对 $D$ 作高斯积分，代数上其实更简单一些：那样甚至不必约化到 $D$ 的零模。但我们认为按上面的次序作积分在教学上更有用。}
% [Tutor Note: $\Theta$ 的结构值得逐项核对：$e^{\frac{\Sigma_+^{2}}{8\pi y}(\frac{\dd u}{\dd a})^{2}}$ 与 $\frac{1}{\sqrt{y}}$ 来自对 $D$ 零模的高斯积分（$D$ 是自对偶 $2$-形式，$b_2^{+}=1$ 时其零模即 $D_+\omega$，配平方后留下常数项与测度因子）；方括号 $\left[\lambda_+ + \frac{\ii}{4\pi y}\frac{\dd u}{\dd a}\Sigma_+\right]$ 是配平方后剩下的“平移”，它使 $\Theta$ 成为 $\lambda$ 上的 theta 函数而非普通高斯和。$y$ 只以 $\lambda_+$ 的系数形式出现，这决定了 (7.12) 之后按 $1/y$ 展开时每一项都是模形式。]

关于 $u$ 平面积分的这一结果，我们作几点评论。

\begin{enumerate}
\item 这一表达式写成对包含超多重态（秩一情形）也成立的形式。后面会派上用场。
\item 我们可以把积分改写为
\begin{equation}
\int\dd a\,\dd\bar a\,(\cdots)=\int_{\CC}\dd u\,\dd\bar u\left|\frac{\dd a}{\dd u}\right|^{2}(\cdots)
\tag{7.8}
\end{equation}
% [Tutor Note: (7.8) 的被积函数大量地（而且是本质地）使用某个对偶标架：$\lambda$ 的 $\pm$ 分解、$\Theta$ 中的 $\frac{\dd u}{\dd a}$ 都依赖所选标架。作者说“$\Theta$ 本质上是 theta 函数、具有良好的对偶变换性质”正是说：在改变对偶标架（$SL(2,\ZZ)$ 的单值性变换）下，测度 $\dd u\,\dd\bar u\,|\dd a/\dd u|^{2}\Theta$ 的变换因子彼此相消，故它确实是 $u$ 平面上的单值测度。检验须分两步：先验证单值性（与标架无关），再验证积分本身的收敛性（见下文对 $u=\infty$ 的分析）。]
然而要注意，被积函数大量地（而且是本质地）使用了某个对偶标架。这一表达式是否真的是 $u$ 平面上良定义的测度，并非显然；但可以借助各项在对偶标架变换下的性质加以检验。在这一讨论中，重要的是 $\Theta$ 本质上是 theta 函数，因而具有漂亮的对偶变换性质。一旦验证了测度是单值的，接下来还必须检验积分确实良定义。这一点颇为微妙，下面将予讨论。
\item 在纯 $SU(2)$ 理论中，更好的写法是
\begin{equation}
\int\dd a\,\dd\bar a\,(\cdots)=\int_{\cF}\dd\tau\,\dd\bar\tau\left|\frac{\dd a}{\dd\tau}\right|^{2}(\cdots)
\tag{7.9}
\end{equation}
其中积分区域 $\cF$ 与 $\Gamma^{0}(4)$ 在上半平面上的基本区域同构：
% [Tutor Note: 从 $a$ 平面换到 $\tau$ 平面：纯 $SU(2)$ 理论中 $a$ 是 $\tau$ 的单值函数，故 $\dd a\wedge\dd\bar a=|\dd a/\dd\tau|^{2}\,\dd\tau\wedge\dd\bar\tau$，雅可比即 $|\dd a/\dd\tau|^{2}$。区域 $\cF$ 是 $\Gamma^{0}(4)$ 的基本区域（原文写作 $\Gamma^{0}(4)$，上标为零），它有三个尖点 $\tau=\ii\infty,\,0,\,2$，这正解释了作者紧接着说“$q$ 展开在三个尖点附近形式各不相同”，也是 (7.12) 之后必须分区域讨论收敛性的原因。注意 $b_2^{+}=1$ 时 $\cF$ 是二维实区域，而 (7.8) 中的 $u$ 平面同构于它——这正是“$u$ 平面积分”这个名字的由来。]
\begin{figure}[H]
  \centering
  \figplaceholder{〔原文此处（印刷页第 61 页，PDF 第 62 页）留有图位“FIGURE OF FUNDAMENTAL DOMAIN FOR $\Gamma^{0}(4)$”，PDF 中未插入实际图形〕}
  \caption{原文图题“FIGURE OF FUNDAMENTAL DOMAIN FOR $\Gamma^{0}(4)$”：$\Gamma^{0}(4)$ 在上半平面上的基本区域，即 (7.9) 中的积分区域 $\cF$。}
  \label{fig:gamma0-4}
\end{figure}
% [原 PDF 第 62 页]

于是被积函数中的全部因子都可写成 $\tau$ 的模函数，尽管在 $\tau=\ii\infty$、$\tau=0$ 与 $\tau=2$ 附近相应的 $q$ 展开形式各不相同。
\item 在判别式轨迹附近以及 $u=\infty$ 附近，被积函数中的若干项会变奇异。在这些区域必须小心地定义积分。作为例子，让我们详细考察 $u=\infty$ 处被积函数的性质。我们的积分可写成
\begin{equation}
Z_u^{\xi}(p,\Sigma)=\int_{\cF}\dd\tau\,\dd\bar\tau\,H(\tau)\Theta
\tag{7.10}
\end{equation}
其中 $H(\tau)$ 只是 $\tau$ 的函数，不含 $\bar\tau$：
\begin{equation}
H(\tau)=\alpha^{\chi}\beta^{\sigma}\frac{\dd a}{\dd\tau}
\left(\frac{\dd u}{\dd a}\right)^{\chi/2}
\Delta^{\sigma/8}
e^{\,2pu+\Sigma^{2}T_c(u)}
\tag{7.11}
\end{equation}

现在考察区域 $\operatorname{Im}\tau\to\infty$。由 (5.41)、(5.42) 这类表达式可得大 $y$ 行为
\begin{align}
u &= \frac{1}{8}q^{-1/4}\left(1+\mathcal{S}(q^{1/2})\right) \notag\\
\frac{\dd a}{\dd u} &= q^{1/8}\left(1+\mathcal{S}(q^{1/2})\right) \notag\\
\frac{\dd u}{\dd a} &= q^{-1/8}\left(1+\mathcal{S}(q^{1/2})\right) \tag{7.12}\\
\frac{\dd a}{\dd\tau} &= \frac{\dd a}{\dd u}\frac{\dd u}{\dd\tau}
= -\frac{2\pi\ii}{32}q^{-1/8}\left(1+\mathcal{S}(q^{1/2})\right) \notag\\
T_c(u) &= q^{1/4}\left(1+\mathcal{S}(q^{1/2})\right) \notag
\end{align}
% [Tutor Note: (7.12) 实为五行（$u$、$\dd a/\dd u$、$\dd u/\dd a$、$\dd a/\dd\tau$、$T_c(u)$），而紧接的正文说“上面四行中的具体级数互不相同”，应是原稿笔误，实指五行。另可核对第三、四行的自洽性：由 $u=\frac18q^{-1/4}(1+\cdots)$ 得 $\frac{\dd u}{\dd\tau}=-\frac{2\pi\ii}{4}\cdot\frac18q^{-1/4}(1+\cdots)=-\frac{2\pi\ii}{32}q^{-1/4}(1+\cdots)$，乘以 $\frac{\dd a}{\dd u}=q^{1/8}(1+\cdots)$ 即得 $\frac{\dd a}{\dd\tau}=-\frac{2\pi\ii}{32}q^{-1/8}(1+\cdots)=-\frac{\pi\ii}{16}q^{-1/8}(1+\cdots)$。$q^{1/8}$ 的出现正是 $\tau\to\tau+4$ 不变性检验的核心（见脚注 34）。]
这里的记号含义如下：$\mathcal{S}(q^{1/2})$ 表示以 $q^{1/2}$ 的正幂次为项的级数，系数为整数，首项为 $q^{1/2}$ 阶。上面四行中出现的具体级数彼此不同——记号 $\mathcal{S}(q^{1/2})$ 只是表示某个这样的级数的通称。回忆 $q=\ee^{2\pi\ii\tau}=\ee^{2\pi\ii x-2\pi y}$，因此我们已把 $y\to\infty$ 时指数增长的那些项与指数衰减的项分离开来。$H(\tau)$ 与 $\Theta$ 在 $y\to\infty$ 处都有奇异性。

净结果是积分具有如下形式：
\begin{equation}
\int_{\cF}\frac{\dd x\,\dd y}{y^{1/2}}
\sum_{\lambda}
e^{-\ii\pi\bar\tau\lambda_+^{2}-\ii\pi\tau\lambda_-^{2}}
(-1)^{(\lambda-\lambda_0)\cdot w_2(X)}
\sum_{\mu\in\frac{1}{8}\ZZ}
q^{\mu}\left(c_{\lambda}^{(0)}(\mu)+y^{-1}c_{\lambda}^{(1)}(\mu)+\cdots\right)
\tag{7.13}
\end{equation}
% [Tutor Note: 关于 $q^{1/8}$ 的幂次与 $\tau\to\tau+4$ 不变性：级数中出现 $q^\mu$、$\mu\in\frac18\ZZ$，即既有 $q^{1/8}$ 的偶次幂也有奇次幂，故单看被积函数它并不明显在 $\tau\to\tau+4$ 下不变。完整检验见脚注 34：$\tau\to\tau+4$ 使 theta 函数中 $\lambda$ 的二次项获得整体相位 $\ee^{-4\pi\ii\lambda^{2}}=\ee^{-\ii\pi w_2(E)^{2}}$，同时 $q^{1/8}\to-q^{1/8}$ 导致 $\frac{\dd u}{\dd a}\to-\frac{\dd u}{\dd a}$（对 $0$-可观测量耦合不变、对 $2$-可观测量耦合变号），须配合 $\lambda\to-\lambda$ 重排求和与测度中 $\frac{\dd a}{\dd\tau}$、$(\frac{\dd u}{\dd a})^{\chi/2}$、$\Delta^{\sigma/8}$ 的相位，最终净相位为 $\ee^{-\ii\pi w_2(E)^{2}}\ee^{-\ii\pi(2\lambda)\cdot w_2(X)}$，由 $\chi+\sigma=4$ 与 $w_2(X)$ 为特征向量知其为 $1$。]
关于这一表达式有若干评论。首先，尽管 $q^{\mu}$ 级数中既有 $q^{1/8}$ 的偶次幂又有奇次幂，但稍作计算即可检验：被积函数在 $\tau\to\tau+4$ 下不变。\footnote{检验此点的细节如下。在 $\tau\to\tau+4$ 下，theta 函数中 $\lambda$ 的二次项整体获得相位 $\ee^{-4\pi\ii\lambda^{2}}$；由于 $\lambda$ 取值于一个挠子，该相位等于 $\ee^{-\ii\pi w_2(E)^{2}}$。其次，$q^{1/8}\to-q^{1/8}$，故 $u\to u$ 而 $\frac{\dd u}{\dd a}\to-\frac{\dd u}{\dd a}$，于是对 $0$-可观测量的耦合不变，对 $2$-可观测量的耦合变号。我们必须通过在 $\lambda$ 的求和中作 $\lambda\to-\lambda$ 来抵消这一变号。这带来额外相位 $-\ee^{-\ii\pi(2\lambda)\cdot w_2(X)}$，其中已计入插入本身给出的符号。与此同时，测度中 $\frac{\dd a}{\dd\tau}\to-\frac{\dd a}{\dd\tau}$，而 $\left(\frac{\dd u}{\dd a}\right)^{\chi/2}\to\ee^{-\ii\pi\frac{\chi}{2}}\left(\frac{\dd u}{\dd a}\right)^{\chi/2}$、$\Delta^{\sigma/8}\to\ee^{-\ii\pi\frac{\sigma}{2}}\Delta^{\sigma/8}$。利用 $\chi+\sigma=4$ 可知测度 $\dd\tau\,\dd\bar\tau\,\frac{\dd a}{\dd\tau}\left(\frac{\dd u}{\dd a}\right)^{\chi/2}\Delta^{\sigma/8}$ 改变一个负号。于是我们得到的净相位为 $\ee^{-\ii\pi w_2(E)^{2}}\ee^{-\ii\pi(2\lambda)\cdot w_2(X)}$；再次利用 $\lambda$ 的量子化条件，这等于 $\ee^{-\ii\pi w_2(E)^{2}-\ii\pi w_2(E)\cdot w_2(X)}$。由于 $w_2(X)$ 是 $H^{2}(X;\ZZ)$ 的特征向量，该相位等于 $1$。}

其次，这一级数相当形式化：由于指数中与 $p$ 和 $\Sigma$ 的耦合涉及 $q^{1/8}$ 的负幂次，当我们把表达式按 $p$ 与 $\Sigma$ 的幂级数展开时，会遇到任意大的 $q^{-1/8}$ 正幂次。若按 $S$ 中次数为 $r$ 的多项式与 $p$ 中次数为 $\ell$ 的多项式来组织被积函数（$r,\ell\ge0$），那么同样有 $1/y$ 的展开，而且对所有 $n\ge0$，$1/y^{n}$ 的系数是一个模表达式，其 $q$ 展开具有如下形式
\begin{equation}
q^{-\frac{3+r+2\ell}{8}}\left(1+\mathcal{S}(q^{1/2})\right)
\tag{7.14}
\end{equation}
（此时级数 $\mathcal{S}(q^{1/2})$ 的系数是复数）。因此一个瞬子模空间 $\cM$ 给出首项发散 $q^{-(6+\dim\cM)/16}$ 加上更高阶的项。

此外，对固定的次数 $r$，$1/y^{n}$ 中的级数在 $n=r$ 处截断。若像 (7.13) 那样固定 $q^{\mu}$ 的一个幂次，则每个 $c_{\lambda}^{(n)}$ 都是 $\Sigma$ 与 $p$ 的级数，且 $r+2\ell\ge-3-8\mu$。$1/y$ 展开中的首项为 $c_{\lambda}^{(0)}(\mu)=\lambda_+c(\mu)$，其中 $c(\mu)$ 是 $p/4,\Sigma$ 的级数，其系数是 $\Gamma^{0}(4)$ 的模形式的系数，取整数值的傅里叶系数。

现在考虑 $y\to\infty$ 时尖点附近积分的定义。我们先对 $x$ 在 $0\le x\le1$ 上积分，得到 $y$ 的一个函数。这一步把对 $\lambda$ 的求和投影，使得只有第零个傅里叶系数有贡献：
\begin{equation}
2\mu=\lambda^{2}
\tag{7.15}
\end{equation}
于是剩下
\begin{equation}
\int_{y_{\mathrm{cutoff}}}^{\infty}\frac{\dd y}{y^{1/2}}
\sum_{\lambda}
e^{-2\pi y\lambda_+^{2}}
(-1)^{(\lambda-\lambda_0)\cdot w_2(X)}
\left(
\lambda_+\,c\!\left(\tfrac{1}{2}\lambda^{2}\right)
+y^{-1}c_{\lambda}^{(1)}\!\left(\tfrac{1}{2}\lambda^{2}\right)+\cdots
\right)
\tag{7.16}
\end{equation}
现在，只要不出现 $\lambda_+=0$ 的项，这一积分绝对收敛。物理上，$\lambda_+$ 所对应的解正是阿贝尔瞬子。阿贝尔瞬子带来额外的玻色零模，而出现额外玻色零模时理应预期新的发散。此外，在 $1/y$ 的级数中，即使 $\lambda_+=0$，也只有首项能导致积分发散。这一潜在的发散正是下文所讨论的穿墙的来源。
% [Tutor Note: (7.16) 对 $\lambda_+$ 的符号依赖是本节的枢纽：$\int_{Y}^{\infty}\frac{\dd y}{y^{1/2}}\ee^{-2\pi\lambda_+^{2}y}\lambda_+$ 在 $\lambda_+\to0^{\pm}$ 时的极限与 $\mathrm{sign}(\lambda_+)$ 成正比（作代换 $t=2\pi\lambda_+^{2}y$ 可见），故 (7.21) 中出现 $\mathrm{sign}(\lambda_+)$。当 $\omega\cdot\lambda\neq0$ 时，对 $\lambda$ 与 $-\lambda$ 两项的这个符号相反、彼此相消，度规依赖为零；只有 $\omega\cdot\lambda=0$ 时 $\lambda_+$ 过零、符号发生跳变，才留下非零的 $\delta$ 函数型贡献。这就是 (7.19) 支撑在墙 $W(\lambda)$ 上的实质。]
\item 尽管这一积分微妙而复杂，我们必须强调：$X$ 的拓扑只通过经典上同调环进入，因此 $Z_{IR,Coul}$ 只是 $X$ 的同伦型的函数。
\item 注意，尽管我们讨论的是拓扑场论，被积函数却确实具有非平凡的度规依赖，因为它显式使用了 $\lambda$ 到其自对偶部分 $\lambda_+$ 与反自对偶部分 $\lambda_-$ 的投影。既然处理的是拓扑场论，我们或许可以期望积分结果与度规无关。下面我们就转向对这一问题的详细研究。
\end{enumerate}

\subsection{度规依赖：穿墙}

拓扑场论的形式体系保证：路径积分对度规的变分将是场空间中的全导数：
\begin{equation}
\frac{\delta}{\delta g^{\mu\nu}}Z=\langle T_{\mu\nu}\rangle=\langle\{Q,\Lambda_{\mu\nu}\}\rangle
\tag{7.17}
\end{equation}
然而在某些情形下这一全导数不为零。BCOV 的全纯反常即是一例。$u$ 平面积分是这一现象的另一个突出例子。

$u$ 平面积分的被积函数完全通过诸如
\begin{equation}
\lambda=\lambda_+\omega+\lambda_-
\tag{7.18}
\end{equation}
这样的投影来获得度规依赖。因此我们只需考虑沿双曲面的一族 $\omega(t)$ 并对 $t$ 求导。可以算出显式的全导数，把 $Z$ 对度规的变分化约成一个边界积分：
\begin{equation}
\frac{\dd}{\dd t}Z_u^{\xi}(p,\Sigma)
=
-\ii\alpha^{\chi}\beta^{\sigma}
\sum_{u_{*}=\pm\Lambda^{2},\infty}
\lim_{\epsilon\to0}
\oint_{S^{1}(\epsilon)}
\dd u\,
\left(\frac{\dd a}{\dd u}\right)^{1-\frac{1}{2}\chi}
\Delta^{\sigma/8}
e^{\,2pu+\Sigma^{2}T(u)}
\Upsilon
\tag{7.19}
\end{equation}
其中 $\Upsilon$ 是另一个与 $\Theta$ 类似的 theta 函数。仔细的分析表明，这是一个 $\delta$ 函数（$N_f=4$ 的情形除外）。该 $\delta$ 函数的支撑位于形如
\begin{equation}
W(\lambda):=\{\omega:\omega\cdot\lambda=0\}
\tag{7.20}
\end{equation}
的某些墙上。
% [Tutor Note: 为什么度规变分化成边界项：$T_{\mu\nu}=\{Q,\Lambda_{\mu\nu}\}$ 是 $Q$-恰当的，故 $\delta Z/\delta g^{\mu\nu}$ 在模空间“内部”的贡献为零（$Q$-恰当量的期望值在 $Q$-闭插入下为零），只剩下 $\cM(\lambda)$ 的“端”处的贡献；在 $\tau$ 平面上这些端就是 $u\to\pm\Lambda^{2}$（两个尖点 $\tau=0,2$ 附近）与 $u\to\infty$，故 (7.19) 写成围绕这些点的小圆 $S^{1}(\epsilon)$ 上的围道积分。为什么是 $\delta$ 函数且支撑在 $W(\lambda)$ 上：当 $\omega\cdot\lambda\neq0$ 时，$\lambda\to-\lambda$ 两项的 $\mathrm{sign}(\lambda_+)$（见 (7.21)）相反相消，变分为零；只有 $\omega\cdot\lambda=0$ 时 $\lambda_+$ 改变符号，$\mathrm{sign}$ 跳变，给出 $\delta$ 函数型的跳变。于是“度规依赖”并不是连续的，而是沿墙的跳变——这正是穿墙现象的几何含义。]
为了更详细地理解这一穿墙现象，我们回到 (7.16) 中对积分的正则化。在 $u\to\infty$ 附近，规范耦合 $\operatorname{Im}\tau\to\infty$（在几乎典范的对偶基下），这会使 theta 函数中除一项之外的所有项指数式压低，留下的是与某个使 $\lambda_+\to0$ 的向量 $\lambda$ 相伴的项。由 (7.16) 可见积分的行为如
\begin{equation}
c\!\left(\tfrac{1}{2}\lambda^{2}\right)
\int^{\infty}\frac{\dd y}{y^{1/2}}\,e^{-2\pi\lambda_+^{2}y}\,\lambda_+
\sim
c\!\left(\tfrac{1}{2}\lambda^{2}\right)\operatorname{sign}(\lambda_+)
\tag{7.21}
\end{equation}
其中 $c(\mu)$ 是 $\wp$ 与 $\Sigma$ 的形式级数，其任何一项的系数是 $\Gamma^{0}(4)$ 的某个模形式的傅里叶系数。
% [Tutor Note: 原文此处用花体 $\wp$ 表示 $0$-可观测量 $\cO(\wp)$ 的插入点（其上同调类即点类，记号与 §5 中 $\ee^{p\cO(\wp)}$ 一致），故 $c(\mu)$ 是点类 $\wp$ 与曲面类 $\Sigma$ 的形式级数；不要与 Weierstrass 椭圆函数 $\wp$ 混淆。另注意 (7.21) 中 $\lambda_+$ 被放在积分号内、$y$ 积分之外：它的符号决定整个积分的符号，故 $\lambda_+\to-\lambda_+$ 时结果变号，这是“墙两侧相消”的技术根源。]
物理上，在 $u=\infty$ 的墙上发生的事情是：联络可以变为可约的，并且存在阿贝尔瞬子，即线丛 $L$ 上满足 $F^{+}=0$ 的联络。这在路径积分中带来额外的玻色零模，从而导致 $\delta$ 函数型的发散。

这些墙分别位于：来自 $u=\infty$ 的墙为
\begin{equation}
\lambda\in\Gamma_{\xi}=\frac{1}{2}\bar{w}_2(P)+\bar{H}^{2}(X)
\tag{7.22}
\end{equation}
来自 $u=\pm\Lambda^{2}$ 的墙为
\begin{equation}
\lambda\in\Gamma_{w}=\frac{1}{2}\bar{w}_2(X)+\bar{H}^{2}(X)
\tag{7.23}
\end{equation}
% [Tutor Note: (7.22) 与 (7.23) 的区别在于“特征类”：$u=\infty$ 处的可约联络给出的是 $\frac12\bar w_2(P)$ 的挠子（$P$ 是 $w_2$ 的提升所对应的主丛），而 $u=\pm\Lambda^{2}$ 处的可约解给出的是 $\frac12\bar w_2(X)$ 的挠子。两族墙的来源不同，故第 8 节必须论证：来自 $u=\pm\Lambda^{2}$ 的跳变在完整的唐纳森配分函数中不出现，只能被希格斯分支的跳变抵消。]
跨过墙的间断 $\Delta_{u_{*},\lambda}Z_u^{\xi}(p,\Sigma)$ 可以表为一个全纯对象的留数：它就是某个模形式的傅里叶系数。

这些墙把前向光锥分割成若干腔室。对固定的 $\ell,r$，关联函数 $\langle\cO^{\ell}\cO(\Sigma)^{r}\rangle$ 只会在有限多个腔室的边界上改变。

\begin{figure}[H]
  \centering
  \figplaceholder{〔原文此处（印刷页第 64 页，PDF 第 65 页）留有图位“FIGURE: CHAMBERS”，PDF 中未插入实际图形〕}
  \caption{原文图题“CHAMBERS”：前向光锥被各道墙 $W(\lambda)=\{\omega:\omega\cdot\lambda=0\}$ 分割所成的腔室。}
  \label{fig:chambers}
\end{figure}
% [原 PDF 第 65 页]

于是任何关联函数的度规依赖都是分片常数的。跨墙的\emph{穿墙公式}将涉及模形式的傅里叶系数。

对来自 $u=\infty$ 的墙，其穿墙公式（WCF）$\Delta_{\infty,\lambda}Z_u^{\xi}(p,\Sigma)$ 精确复现了 L.~格策 (格策) 关于 $b_1=0$、$b_2^{+}=1$ 时唐纳森多项式变化的公式，只要取
\begin{equation}
\alpha^{\chi}\beta^{\sigma}=\frac{2^{(2+3\sigma)/4}}{\pi}
\tag{7.24}
\end{equation}
% [Tutor Note: (7.24) 只定出组合 $\alpha^{\chi}\beta^{\sigma}$。由 $\chi+\sigma=4$ 可把它改写为 $\alpha^{4}(\beta/\alpha)^{\sigma}$，故 (7.24) 在 $\sigma$ 变化下的相容性给出 $\alpha^{4}=2^{1/2}/\pi$、$\beta/\alpha=2^{3/4}$，即 $\alpha=2^{1/8}\pi^{-1/4}$、$\beta=2^{7/8}\pi^{-1/4}$；因 $\chi$ 为偶数，$\alpha$ 的符号不影响结果。原文随后说“这完全固定了 $\alpha,\beta$”，正是指这一分解。]

现在 $\chi+\sigma=4$，而 $\sigma=1-b_2^{-}$ 可以变化，故这完全固定了 $\alpha,\beta$。（在这些方程中我们还把 $\Lambda$ 标度为 $1$。）

然而，我们还从 $u=\pm\Lambda^{2}$ 处的奇异性得到跨墙的穿墙公式。既然我们已经用 $\Delta_{\infty,\lambda}Z_u^{\xi}(p,\Sigma)$ 完全解释了唐纳森多项式的变化，这些新的间断就\emph{不能}是完整配分函数 $Z_{DW}^{\xi}(p,\Sigma)$ 的间断。
